\documentclass[11pt,a4paper]{article}

\usepackage[margin=2.5cm]{geometry}
\usepackage{amsmath,amssymb,amsfonts}
\usepackage{authblk}
\usepackage{hyperref}
\usepackage{enumitem}
\setlist[itemize,1]{label=---}

\title{Proper Time as Coherence-Cycle Counting and Emergent Metric in a Relational Graph Framework (PDL)}

\author{Cédric Laubscher}
\affil{Independent researcher, Switzerland\\
\texttt{[cedric.laubscher@gmail.com]}}

\date{}

\begin{document}

\maketitle

\begin{abstract}
We introduce the Projective Dynamic Logo (PDL), a relational, graph-theoretic framework that reconstructs proper time, metric structure, and gravitation from purely combinatorial axioms formulated on finite signed graphs. Within this setting, physical entities are modeled as minimal closed configurations of binary agreement/disagreement relations.

A foundational optimization principle—comprising triangular coherence together with configurational minimality—selects, as the first non-trivial stationary structure, a complete signed graph on four vertices with six edges. Its internal 2-cycle is interpreted as a logical prototype of a rest-frame electron, with proper time identified with discrete counts of such coherence cycles.

Composite structures (such as nucleons) are represented as hierarchical assemblies of locally stationary regimes embedded in a dynamical “relational sea.” An emergent metric structure arises from compatibility relations among simultaneously existing closures, whereas gravitation is modeled as spatial variation in a “coherence cost” functional, induced by a small but systematic leakage of coherence from composites into their environment.

This framework thus provides a causal, relational account of proper time and spacetime in the absence of a background manifold, and yields structurally constrained expressions for physical constants—such as the fine-structure constant—interpreted as coherence ratios within the same discrete combinatorial architecture.

\end{abstract}

\noindent\textbf{Keywords:} causal structure; proper time; emergent spacetime; relational framework; graph-based quantum gravity; discrete signed graphs.

\section{Motivation and scope}

Clarifying the role of causality, temporal ordering, and spacetime structure in quantum and relativistic physics remains a central open challenge. Standard formalisms typically assume from the outset a continuous spacetime manifold and subsequently analyze causal relationships within this fixed geometric background. By contrast, \emph{Projective Dynamic Logo} (PDL) adopts a fundamentally relational perspective, in which the primitive entities are finite signed graphs, selected according to coherence and minimality principles. Within this setting, one investigates the conditions under which quantities such as proper time, metric structure, and gravitational interaction can arise as emergent, derived notions.

The objective of this extended abstract is to outline how, within the PDL framework, (i) proper time can be reinterpreted as a discrete counting of coherence cycles associated with minimal graph configurations, and (ii) an effective metric and gravitational coupling can be induced from compatibility conditions and variations of a coherence cost functional over a network of such closures. We also briefly indicate how this formulation interfaces with questions of causal structure and emergent spacetime that are central to the Causalworlds program.

\section{Relational axioms and minimal stationary structure}

Within the framework of PDL, a \emph{logical configuration} is modeled as a finite signed graph
\[
G = (V,E,s),
\]
where $V$ is a finite set of vertices (representing informational entities), $E \subseteq V \times V$ is a set of undirected edges, and $s : E \to \{-1,+1\}$ assigns to each edge a binary sign encoding agreement ($+1$) or disagreement ($-1$). Triangles (3-cycles) are treated as the elementary motifs of logical closure: for any $i,j,k \in V$ such that $(i,j)$, $(j,k)$, and $(i,k)$ are edges of $G$, the triangle is defined to be \emph{coherent} if
\[
s_{ij}\, s_{jk}\, s_{ik} = +1,
\]
and \emph{violated} otherwise. The degree of incoherence is quantified by a coherence-defect parameter $\varepsilon$, specified as the proportion of violated triangles among all unordered triples of distinct vertices.

The construction is constrained by four axioms: minimal binary pulsation, triangular coherence, minimal completeness, and logical optimisation. Collectively, these axioms select configurations that (i) admit a non-trivial two-step update of edge signs (a binary pulsation) that leaves coherence invariant, (ii) minimise the coherence defect $\varepsilon$, (iii) implement logical closure without recourse to any external reference structure, and (iv) are minimal with respect to set-theoretic inclusion of vertices and edges.

Subject to these axioms, the smallest non-trivial elementary stationary structure is a complete signed graph on four vertices and six edges, denoted $(4,6)$. There exists an assignment of signs for which all triangular subgraphs are coherent ($\varepsilon = 0$), the set of coherent triangles forms a connected covering of the graph, and a global sign inversion induces a 2-cycle $s \mapsto -s \mapsto s$ that preserves coherence at every step. This $(4,6)$ configuration is interpreted as a logically stationary regime and will be used as a prototype for an electron in a rest state.

\section{Proper time as coherence-cycle counting}

The intrinsic 2-cycle of the $(4,6)$ structure specifies a minimal logical alternation that can be iterated indefinitely without degradation of coherence and without reference to any \emph{a priori} temporal metric. Within the PDL framework, the proper time of an elementary entity is reinterpreted as the discrete counting of these internal coherence cycles.

More precisely, if the $(4,6)$ structure is associated with an electron in its rest frame, its internal pulsation can be identified with the Compton frequency, and a single completed cycle can be associated with a minimal quantum of action. Proper time along a worldline is thereby represented as the number of internal coherence cycles sustained by the corresponding closed structure, rather than as a continuously valued parameter introduced externally. Relativistic time dilation can then be reformulated as a modification of the coherence budget required to maintain the internal dynamical regime across distinct relational environments.

In this perspective, proper time is not a primitive coordinate but a discretized counting of logically constrained alternations supported by closed graph configurations. This provides a route toward reformulating causal order and temporal structure in terms of the relations among such coherence cycles.

\section{Composite structures, emergent metric, and coherence cost}

Beyond the elementary $(4,6)$ closure, PDL accommodates composite entities, such as nucleons, as hierarchical assemblages of local stationary regimes embedded within a dynamically evolving “relational sea.” A concrete proton realisation is modelled by three local stationary domains (valence cores), supplemented by an extensive set of dynamical relations specified by integer parameters that encode the numbers of stationary and dynamical edges. This construction yields a closed but internally structured graph endowed with a finite “active surface,” through which it can interact with an external $(4,6)$ electron-type structure.

In this framework, energy is reinterpreted as the coherence cost required to maintain a given stationary regime, and an effective metric structure is derived from compatibility conditions among simultaneously realised closures. Gravitation is then conceptualised as emerging from spatial and temporal variations in the coherence cost within this induced metric, driven by a small but systematic leakage of coherence from composite structures into their surrounding relational environment. At a coarse-grained level, these variations can be mapped to an effective gravitational coupling and curvature, without positing a fundamental gravitational field on a pre-specified manifold.

\section{Causal structure and emergent spacetime}

The PDL framework motivates a perspective in which causal structure and spacetime geometry are not taken as fundamental, but instead arise from a more primitive relational substratum. Proper time is associated with internal coherence cycles, while causal compatibility between distinct dynamical regimes is characterised by the possibility of sustaining their respective internal cycles without generating mutual inconsistency. In this setting, the emergent metric functions as an effective summary of the “coherence-cost” landscape associated with families of closures, and gravitational phenomena are interpreted as corresponding to gradients within this landscape.

This viewpoint establishes several points of contact with research programmes that regard causal order and spacetime as emergent, as well as with approaches that investigate cyclic and indefinite causal structures. Although PDL is formulated in a discrete, graph-theoretic language rather than in terms of quantum channels or process matrices, it gives rise to analogous foundational questions: how can causal relations and temporal symmetries be defined at a purely relational level, and by what mechanisms can familiar spacetime structures be recovered as effective, coarse-grained descriptions?

\section{Conclusion and outlook}

We have outlined how the PDL framework reinterprets proper time in terms of coherence-cycle counting associated with a minimal $(4,6)$ structure, and how an emergent spacetime metric together with an effective gravitational coupling can be derived from compatibility relations and variations in coherence costs within a hierarchical system of closures. In this way, PDL yields a relational, graph-theoretic account of temporal ordering, causal structure, and spacetime geometry.

A detailed technical exposition, including the combinatorial construction of composite structures and structural expressions for constants such as the fine-structure constant, is provided in:
\begin{itemize}
    \item C.\ Laubscher, \emph{The Emergence of Physical Reality in Projective Dynamic Logo (PDL)}, Zenodo (2026), \href{https://doi.org/10.5281/zenodo.18462686}{doi:10.5281/zenodo.18462686}.
\end{itemize}

Future research will aim to clarify the relationship between PDL and established causal frameworks, for instance by investigating how PDL closures can be mapped to effective process-theoretic descriptions, and by exploring how cyclic or indefinite causal structures may be represented and analyzed within this formalism.

\begin{thebibliography}{9}

\bibitem{LaubscherPDL}
C.~Laubscher,
\newblock \emph{The Emergence of Physical Reality in Projective Dynamic Logo (PDL)},
\newblock Zenodo (2026),
\newblock \href{https://doi.org/10.5281/zenodo.18462686}{doi:10.5281/zenodo.18462686}.

\bibitem{LaubscherIntro}
C.~Laubscher,
\newblock \emph{Introduction to the Projective Dynamic Logo (PDL)},
\newblock Zenodo (2026),
\newblock \href{https://doi.org/10.5281/zenodo.18463130}{doi:10.5281/zenodo.18463130}.

\bibitem{LaubscherBorn}
C.~Laubscher,
\newblock \emph{Relational Emergence of Born's Rule and a Golden-Ratio Active Surface in PDL},
\newblock Zenodo (2026),
\newblock \href{https://doi.org/10.5281/zenodo.18509648}{doi:10.5281/zenodo.18509648}.

\bibitem{LaubscherAlpha}
C.~Laubscher,
\newblock \emph{A Structural Derivation of the Fine-Structure Constant from the PDL Framework},
\newblock Zenodo (2026),
\newblock \href{https://doi.org/10.5281/zenodo.18580925}{doi:10.5281/zenodo.18580925}.

\bibitem{LaubscherPosition}
C.~Laubscher,
\newblock \emph{PDL: Status, Open Problems, and Interface with Other Ontological Frameworks},
\newblock Zenodo (2026),
\newblock \href{https://doi.org/10.5281/zenodo.18581453}{doi:10.5281/zenodo.18581453}.

\end{thebibliography}

\end{document}
